\chapter[并非一时兴起]{并非一时兴起}

\chaptersubtitle{《那真是太令人高兴了》}

\chapternote{%
    ``You do not become a `dissident' \textbf{just because you decide one day to take up this most unusual career}.
    You are thrown into it by your personal sense of responsibility, combined with a complex set of external circumstances.
    You are cast out of the existing structures and placed in a position of conflict with them.
    It begins as an attempt to do your work well, and ends with being branded an enemy of society.''
    \par
    {\raggedleft\noindent ---Václav Havel, \emph{The Power of the Powerless}\par}
    \par
    “一个人之所以成为一名‘不同政见者’，并非仅仅因为他某天忽然决心投入这项非凡的事业。
    他自身的责任感以及各种复杂的外在因素驱使他加入这一事业。
    他被现行制度排斥，并被置于与之冲突的位置。
    事情始于努力做好本职工作的美好愿望，却以被贴上社会公敌的标签告终。”
    \par
    {\raggedleft\noindent ——瓦茨拉夫·哈维尔《无权者的权力》\par}
    \par
    译者按：成为魔法少女并非兴之所至，“变身”亦是被责任感与复杂的外在境遇推入的身份。
}

魔法少女的变身动画有着耐人寻味的历史，也承担着特殊功能，这使麻美的变身再度成为《魔圆》第二集开场片段的不二之选。
该场景最早起源于 1973 年永井豪\footnote{%
    译注：永井豪（1945 年 9 月 6 日—），本名永井洁，日本漫画家，东京都立板桥高等学校毕业。
    代表作有《恶魔人》《魔神Z》《甜心战士》等。
}的漫画及改编动画《甜心战士》，其以自身兼具魔法少女与少年漫画（即“为男孩”设计）元素为卖点\cite{saito2011beautiful-fighting-girl}。
这很快便成为了魔法少女题材的标志性特点。
到 1990 年代，这类作品或多或少都会融入一些少年漫画元素\cite{orbaugh2003busty-battlin-babes}。

在大多数魔法少女作品中，变身动画发挥着多重作用。
传统场景（以《美少女战士》为代表）通常会展示一连串裸体角色的彩色剪影，而后魔法少女服装在其身上逐渐形成。
魔法少女摆出一系列姿势，镜头围绕她旋转，营造出一个裸体女孩独自在代表男性凝视\footnote{%
    这里“男性凝视”的概念指电影和电视镜头构图中一种倾向，即镜头模仿异性恋男性观众的窥视方式，例如通过在镜头中长时间聚焦于女性身材曲线。
}的镜头\cite{orbaugh2003busty-battlin-babes,mulvey2006visual-pleasure}前舞动的效果。
从叙事理论角度看，这一瞬间象征着角色赋权，伴随背景音乐，它也会唤起观众相应的情感。
然而，镜头的运作同时也将角色塑造成供人观看的对象，换句话说，她虽被赋予力量，但这种力量十分有限。
她依然在展示少女气质，并服从于霸权男性气质\footnote{%
    简而言之，霸权男性气质指的是一种文化倾向，即将男性气质与权力等同。
    暴力、支配力和果断等权力表现被视为固有或适当的男性特征。
    相比之下，表演女性气质指的是一种文化倾向，即将女性气质视为一种表演，一种社会游戏，女性通过这种表演为自我表达创造一个有限但相对安全的空间，但代价是迎合霸权男性气质。
}\cite{connell1987gender-power}。
这样的变身场景基本上是在传达：
“她是一位强大的守护者，但依然不会威胁到父权主义或脆弱的男性自尊。
她的存在目的，不过是为了取悦你。”
这正是魔法少女题材逐渐与刻板印象中的变态“御宅\footnote{%
    例如，\emph{Anime News Network} 上对《魔法少女奈叶》的批评文章，称其为“面向御宅的娱乐，已被推到其逻辑极限”，特别提到其“萌系角色……和色情服务镜头”。
}”产生联系的原因\cite{kimlinger2008nanoha-review}。

日本精神科医生斋藤环\footnote{%
    译注：斋藤环（1961 年 9 月 24 日—），日本精神病学家、评论家，筑波大学医学医疗系社会精神保健学教授。
    专门研究青春期精神病理与精神病志学。
}在其著作《战斗美少女的精神分析》中指出，“御宅”不仅可指“技术迷、书呆子、粉丝”，还带有比这些词更为尖酸刻薄的含义，指那些将所消费的虚构作品视为自身身份的一部分，并从中寻找性欲对象的人。
换句话说，日本御宅文化的一个重要元素（在美国极客文化\footnote{%
    译注：极客（Geek）是指对某一领域有着极大兴趣和热情的人，通常是技术、科幻或游戏等领域。
    该词在美国文化中有时带有贬义，但也被部分人自我认同为一种积极的身份标签。
    这里狭义指代美国的二次元文化。
}中也可以观察到）是作品及其中的角色属于观众，而这种归属感与（通常是异性恋男性的）性欲望息息相关。
正如斋藤环所说：
“变身动画代表了一种加速成熟的过程，实际上是动画给予观众的默许，允许他们对角色抱有性幻想。”
他认为，魔法少女（或者说，他所谓的“战斗美少女”，意指这一概念还可扩展至科幻等其他题材中的角色）并非赋予女性观众强大力量的工具，也不是教育女孩传统女性美德的方式。
她并不是观众想要成为的榜样，而是观众想要“拥有”、通常也会寄托性幻想的对象\cite{saito2011beautiful-fighting-girl}。

相较于大多数同类作品，麻美的变身动画由于并未让她脱去校服，显得更含蓄一些，但镜头仍刻意在胸部和臀部停留，强调身材与服装而非面部特征。
即使麻美在剧中以暴力消灭了威胁小圆的使魔，她仍是一个典型的少女形象，温柔而富有母性。
几分钟后，询子也经历了类似的变身。
这是她在本集中唯一一次出场。
作为一位母亲，她温柔地责备小圆昨晚没有回家。
但她化好妆后，就变成了一位冷静、雄心勃勃的公司高管，盘算着联盟关系与公司内部的权力斗争。
麻美和询子的两个场景可以看作从传统女性角色（学生或母亲）向传统男性角色（战士或胜者）的过渡。

尤其对于麻美而言，变身场景不仅是呈现魔法少女的常规元素，还是一种打破魔女迷宫异化感、强化魔法少女题材与父权社会规范的方式。
这也是一场为小圆与沙耶香而进行的表演，意味着麻美前来拯救她们。
值得注意的是，当麻美在本集结尾再次变身时，并没有采用完整的变身场景；
短暂的光芒与丝带交织过后，变身便已完成。
完整的变身场景强调魔法少女体内存在一个供观众消费的“女学生”形象（小圆和沙耶香处于叙事层内，男性凝视处于叙事层外\footnote{%
    译注：法国文学理论家热拉尔·热奈特，区分了叙事小说中的三种“叙事层次”，分别是：
    外叙事层，指不属于故事本身的叙述者层次，即现实世界的作者；
    叙事层，指故事角色的层次，包括他们的思想和行动；
    元叙事层，指“故事中的故事”，例如故事角色自己讲述的一个故事。
}）。
简短的变身场景则恰恰相反，它提醒观众，少女体内存在“战士”，并不给观众拥有和消费她的机会。
它并非强化社会秩序，反而破坏了秩序，暗示魔法少女并不需要通过为男性观众表演或强调女性气质来获得战斗许可。
她所需要的，仅仅是一个契机。

但对沙耶香而言，这无疑已经足够。
麻美作为无辜者的守护者这一形象立刻吸引了她；
她毫不犹豫地支持麻美，讨厌焰。
对她来说，魔法少女代表着一种明确的善恶对抗。
她迫不及待地想站在正义的一方，渴望成为一名魔法少女，站在麻美身边为正义而战，因而决定带着充作武器的棒球棒参加猎杀魔女的任务。
与此相反，小圆则更在意成为魔法少女的体验。
她想理解麻美和焰，与她们成为朋友。
小圆想成为魔法少女，是因为那样既酷又特别；
她希望自己能像麻美一样，为此还专门设计了一套魔法服装。
换句话说，沙耶香专注于“做”，而小圆则专注于“成为”。
沙耶香果敢行事，却茫然无措；
小圆则目标明确，但踟蹰不前。

尽管两人成为魔法少女的目的有所不同，但都纠结于要许下什么样的愿望。
出乎意料的是，稚嫩的女孩们竟会反思自己所享有的优越生活。
毕竟，作为孩子，她们很难了解自身生活境遇的全貌。
但不得不承认，她们的安全、健康、良好的教育背景、较高的社会地位、出众的外貌以及优渥的经济条件（住在公寓的沙耶香看起来不如小圆富裕；小圆的房子非常大；
但从后续与杏子的互动来看，两人似乎都未曾面临温饱或住房问题），都源于她们有幸出生在特定家庭，而非自身努力。
正如我们稍后将看到的，沙耶香从恭介身上意识到，无常的命运可以轻易打乱一个人的生活。
小圆也认同这个观点，两人都感到自己太过幸运，以至于没有任何愿望值得她们付出丘比所要求的代价。

还有一个隐含的主题贯穿本集：
魔法少女是某种生态系统的一部分。
杏子在几集后明确指出了这一点，但在本集中我们就已经了解到了：
在这个生态系统中，魔女以人类为食，魔法少女则以魔女为食；而丘比是分解者，以捕食行为的副产品为食。
但从社会角度考虑，它却是一个纯粹的捕食者，专门挑选群体中最脆弱的成员，利用她们为自己谋取利益。
从麻美、焰和杏子各自与丘比签约的经历中可以看出，它专门寻找那些正处于痛苦与绝望中的女孩，她们不仅愿意立刻接受契约，而且也更可能迅速转化为魔女。
而小圆与沙耶香迟迟没有接受它的提议，正是因为她们太过幸运；
要操控那些无欲无求的人是格外困难的。

但通常而言，这本不该是魔法少女面临的问题。
丘比寻找愿意为它战斗的女孩，与她们签订契约。
就前两集而言，它的行为看起来比《魔卡少女樱》和《美少女战士》中的对应角色要好些。
在这两部作品中，小樱和小兔都没有选择的余地。
小樱被告知，因为她不小心释放了库洛牌，所以有责任收回这些卡片\cite{cardcaptor-sakura-episode1-2014}；
而小兔则只是被告知，她注定要成为美少女战士。
正如现实生活与虚构作品中常见的那样，女性常常被要求无条件地接受特定角色，无论她们是否同意\cite{connell1987gender-power}。
（这并不是说男性就不会面临类似情况，只是他们更可能被给予选择的机会，并且更容易抗拒自己不得不战斗的命运。例如《新世纪福音战士》中，真嗣不断挣扎着反抗作为 EVA 驾驶员的命运，而小兔只能顺从接受成为魔法少女的命运。）

当然，了解后续剧情之后，丘比狡诈的本质便显而易见了。
它根本不是在寻求知情同意，而是在以令人震惊的方式欺骗和操控少女。
即使在开头的剧情里，有关魔鬼交易和骗子精灵的故事也暗示了丘比并不像表现出来的那样值得信任。
魔女受害者的涂鸦引用了歌德\footnote{%
    译注：约翰·沃尔夫冈·冯·歌德（1749 年 8 月 28 日—1832 年 3 月 22 日），德国戏剧家、诗人、自然科学家、文艺理论家，为古典主义最著名的代表。
}的《浮士德》，进一步强化了这一暗示。
它那张难以捉摸的面孔进一步加剧了这种不安感。
它最初看起来或许可爱友好，但随着观众疑心渐增，那种缺乏情感的表情开始显得像一副面具，让它愈发难以读懂，也进一步加深了观众的猜疑。
小圆和沙耶香也并非生活在《魔卡少女樱》或《美少女战士》那样的世界中（尽管小圆在本集开头醒来的场景与《魔卡少女樱》第二集\footnote{%
    如同《魔圆》里的小圆一样，在《魔卡少女樱》第二集“小樱最好的朋友”的开头，小樱从床上醒来，以为前一集的事情只是一个梦罢了，直到她注意到床边那个根本不是毛绒玩具。
}\cite{cardcaptor-sakura-episode2}有些相似）。
但无论她们是否与丘比签订协议，魔法少女的世界很快也将渗透并侵蚀她们原本安全、舒适、优越的生活。

麻美希望保护无辜者，维系着传统魔法少女题材应有的一切，但她终究会失败。
尽管她将沙耶香的球棒变为传统魔法少女题材中常见的符号——
一根粉红色的卡通魔杖——
但这根魔杖在本集结尾的魔女战斗中并没有派上多少用场，只是帮助沙耶香与小圆暂时脱离险境，而麻美则继续战斗。
同时，\emph{Magia} 这首预示着黑暗版《魔圆》的曲子在战斗中响起，魔女将其超现实的迷宫投射到了小圆的世界里。
尽管麻美再次战胜了魔女，但她在那场精彩战斗中显露出的过度自信（尤其是在她陷入蝴蝶绳索陷阱时），表明她始终在为注视着她的观众表演，而非专注于眼前魔女带来的威胁。

在本集的结尾，麻美已经明确了自己作为沙耶香与小圆的守护者兼引导者的角色定位。
只要她还在，她便会阻止 \emph{Magia} 版《魔圆》，不让其吞没《明天再见》版的魔法世界。
她保护着小圆那个幸福安稳的子供向小世界，使其不被魔女所代表的强烈情绪、复杂性与威胁打破。
然而，这一集也揭示了她的弱点。
显而易见，麻美必将因过度自信而成为牺牲品离开舞台，这样剧情才能继续发展下去。
